\section{Experimentos y validación}

\subsection{Principios de diseño experimental}

Cada DSPy experiment debe dejar en claro `hypothesis`, `metric` y `baseline`. Se recomienda un three-phase pipeline: 1) small-scale data → 2) evidence capture → 3) production dry-run. Compartir el evidence spreadsheet entre equipos permite que otros equipos repliquen el experimento.

\subsection{Simulation to Production}

Experimental validation no puede detenerse en offline benchmark, también se necesitan `sim-to-prod` steps. Cada optimizer change debe ejecutarse con 1000+ queries en el sandbox environment, y registrar `metric drift` y `cost per query`, para garantizar que el rollout sea controlable.

\begin{knowledgebox}{Simulation Pipeline}
El Simulation pipeline incluye `toy dataset` run, `evidence replay`, `sanity check`, y cada paso tiene un defined owner. Si el resultado de evidence replay difiere en más de 5\%, hay que volver a hacer tune o abort.
\end{knowledgebox}

\subsection{Resumen de la sección}

La sección de experimentos y validación ayuda a los equipos a hacer que los ajustes de DSPy sean reproducibles, y genera confianza antes del production rollout.

